\chapter{Introduction}
\label{ch:intro}

\section{Context and Educational Grounding}
Frontier Large Language Models (LLMs) function primarily as ``solution engines.'' They are explicitly optimised, through extensive Reinforcement Learning from Human Feedback (RLHF) and instruction tuning, to deliver immediate, fluent, and direct answers to user queries with minimal friction. While highly effective for productivity and information retrieval, this paradigm represents a fundamental failure mode in educational settings. According to cognitive science theories of learning, particularly Bjork's \textit{desirable difficulties} \cite{bjork1994memory} and Sweller's \textit{Cognitive Load Theory}, learning is not a passive transfer of data, but a generative process. Long-term schema formation and retrieval strength are enhanced when the human brain is forced to engage in active cognitive effort, retrieval practice, and error correction. By providing immediate answers, commercial LLMs bypass this active learning loop, reducing the human student's cognitive engagement to passive consumption.

To address this misalignment, this research presents a Socratic tutoring system designed to act as an active pedagogical guide. The ultimate vision of this work is to deploy this Socratic agent into live student-facing programming environments, enabling personalised, real-time scaffolding. However, designing and optimising the tutor's policy directly with human cohorts is bottlenecked by the high financial, temporal, and pedagogical risks of early-stage evaluations. To overcome this hurdle, we utilise a multi-agent simulation framework as a validation sandbox, simulating student profiles to bootstrap and calibrate the tutor's pedagogical policies in a controlled environment. We utilise a cyclic graph representation of dialogue state in LangGraph to control dialogue flow, update a dynamic learner model, and enforce strict pedagogical boundaries. Once this simulation-based optimisation validates these boundaries, the system is architected to transition seamlessly to live human cohorts via database-backed state interrupts.

\section{Problem Statement: The Competence Paradox}
The direct-answering behaviour of commercial LLMs induces the \textit{Competence Paradox}: a phenomenon where students interact with highly fluent models, copy-paste generated code, and experience a high illusion of competence \cite{prather2024widening}, only to fail independent assessments. From a cognitive perspective, the immediate resolution of programming errors by an LLM prevents the student from building necessary mental models. Actively retrieving information and resolving conceptual friction (cognitive struggle) is necessary to transfer knowledge from working memory to long-term storage. Our system utilises Socratic dialogue to force students to recall programming principles, predict execution outcomes, and explain their reasoning, thereby increasing retrieval strength \cite{roediger2006test}.

Evaluating Socratic tutors using monolithic prompt-engineered student simulators (e.g.~\cite{markel2023gpteach, xu2024eduagent}) is also deeply flawed due to their inability to model cognitive constraints and their susceptibility to conversational alignment. Because standard LLMs are trained to be helpful, polite, and cooperative, simulated student agents exhibit a severe \textit{Sycophancy (Yes-Man) Bias}. When a tutor agent asks, ``Do you understand this concept?'', the simulated student will falsely claim mastery and agree with the tutor's explanations to align with the conversational direction, despite possessing unresolved misconceptions in its underlying state. This makes standard multi-agent educational evaluations unreliable.

\section{Research Questions (RQs)}
We target the intersection of multi-agent state architectures, pedagogical prompting, and automated educational evaluation. Our dissertation is framed around three primary Research Questions:
\begin{itemize}
    \item \textbf{RQ1 (Architectural State Tracking):} \textit{How can a cyclic multi-agent graph with centralised state buffers and a strict single-writer policy dynamically track a learner's epistemic state to inform Socratic pedagogical scaffolding without violating pedagogical constraints (e.g., preventing direct answer leakage)?}
    \item \textbf{RQ2 (Pedagogical Scaffolding \& Efficacy):} \textit{To what extent does a multi-agent Socratic feedback loop, when compared across dyadic (1:1) and multi-party social configurations containing simulated peer agents, improve the normalised learning gain ($g$) of simulated student agents compared to standard direct-answering or linear prompt-based delivery?}
    \item \textbf{RQ3 (Cognitive Realism \& Simulation Validity):} \textit{How closely do simulated student behaviours correlate with human learning characteristics under dynamic cognitive load modelling, and how effectively does a causal inference validation protocol with negative control outcomes mitigate the sycophancy bias and user drift in simulated educational evaluation?}
\end{itemize}

\section{Summary of Technical Contributions}
This research introduces three core innovations to the field of AI in Education (AIED) and Multi-Agent Systems (MAS):
\begin{enumerate}
    \item \textbf{Cognitive Bandwidth \& Friction Modelling (CBFM):} A mathematical formulation representing simulated working memory dynamics based on Cognitive Load Theory, preventing agents from processing overloaded explanations.
    \item \textbf{Sycophancy Mitigation via Adversarial Student Probing:} An information-theoretic bottleneck where student agents must pass independent diagnostic assessments in a private scratchpad to prove genuine conceptual mastery.
    \item \textbf{POMDP Tutor Policy Optimisation:} Formulating Socratic tutoring as a Partially Observable Markov Decision Process (POMDP), allowing systematic, zero-budget comparison of prompt-scaffolding strategies.
\end{enumerate}
